\chapter{Segmentacao, Personas e Digital Twins}
\label{chap:segmentacao_personas_twins}

% ---------------------------------------------------------------
\section{Escopo do capitulo}
% ---------------------------------------------------------------

\begin{quote}\itshape
\textbf{Onde estamos na hist\'oria:} loop interno, intelig\^encia sobre o cliente.
O sistema j\'a responde perguntas gen\'ericas. Agora ele precisa conhecer
quem pergunta: segmento, perfil e hist\'orico. Segmenta\c{c}\~ao, Personas
e Digital Twins transformam o agente gen\'erico em consultor personalizado
--- e preparam a avalia\c{c}\~ao por segmento do cap.~13.
\end{quote}

Este capitulo implementa, em formato de aula pratica guiada, a transformacao
de dados de clientes em tres tipos de agentes de pesquisa: agente por
segmento, agente \textit{persona}\index{persona (IA)|textbf} e
\textit{digital twin}\index{digital twin|textbf}. O fluxo combina
clusterizacao\index{clusterizacao|textbf} com K-Means\index{K-Means|textbf},
geracao de contexto RAG por grupo e criacao sob demanda de agentes para
respostas em primeira ou terceira pessoa conforme o objetivo analitico.

% ---------------------------------------------------------------
\section{Objetivos de aprendizagem}
% ---------------------------------------------------------------

Ao concluir este capitulo, o leitor sera capaz de:

\begin{itemize}[leftmargin=*]
  \item Treinar e aplicar K-Means para segmentacao de clientes;
  \item Construir documentos RAG por cluster e por cliente;
  \item Diferenciar os modos \texttt{segmento}, \texttt{persona} e \texttt{twin};
  \item Validar, com casos controlados, a consistencia de resposta de cada modo.
\end{itemize}

% ---------------------------------------------------------------
\section{Conceitos teoricos}
% ---------------------------------------------------------------

\subsection{Tres superficies de analise}

O projeto adota tres superficies complementares:

\begin{itemize}[leftmargin=*]
  \item \textbf{Segmento:} resposta gerencial em terceira pessoa com base
        no perfil medio do cluster;
  \item \textbf{Persona:} arqu\'etipo nomeado de um cluster, respondendo em
        primeira pessoa representativa;
  \item \textbf{Twin:} simulacao individual, baseada em dados reais do
        cliente e contexto RAG filtrado por \texttt{cliente\_id}.
\end{itemize}

\subsection{Features de clusterizacao}

O vetor de entrada utiliza, minimamente, as variaveis:
\texttt{idade}, \texttt{renda\_mensal}, \texttt{saldo\_medio},
\texttt{transacoes\_mes}, \texttt{score\_credito} e
\texttt{num\_produtos}. Antes do K-Means, aplica-se
padronizacao\index{padronizacao (ML)|textbf} com
\textit{StandardScaler}\index{StandardScaler}.

% ---------------------------------------------------------------
\section{Implementacao orientada por passos}
% ---------------------------------------------------------------

\subsection{Passo 1 --- Gerar base sintetica e executar clustering}

O laboratorio inicia com dados sinteticos para reproducao local:

\begin{lstlisting}[language=Python,
  caption={Treino de clusterizacao e previsao de cluster.},
  label={lst:clustering_lab}]
df = gerar_dados_sinteticos(n=1000)
kmeans, scaler, perfis = executar_clustering(df, n_clusters=4)

cliente_exemplo = {
    "idade": 55,
    "renda_mensal": 18000,
    "saldo_medio": 120000,
    "transacoes_mes": 25,
    "score_credito": 820,
    "num_produtos": 7,
}

cluster_id = int(kmeans.predict(scaler.transform([[ 
    cliente_exemplo["idade"],
    cliente_exemplo["renda_mensal"],
    cliente_exemplo["saldo_medio"],
    cliente_exemplo["transacoes_mes"],
    cliente_exemplo["score_credito"],
    cliente_exemplo["num_produtos"],
]]))[0])
\end{lstlisting}

\subsection{Passo 2 --- Criar contexto RAG por cluster}

Cada cluster gera documentos de contexto com perfil medio, produtos e
estrategia de abordagem. O material e indexado no OpenSearch para
recuperacao posterior.

\begin{lstlisting}[language=Python,
  caption={Geracao de documentos RAG por cluster.},
  label={lst:docs_cluster_lab}]
for cid, perfil in perfis.iterrows():
    docs = _gerar_docs_cluster(cid, perfil)
    indexar_docs_cluster(cid, docs)  # abstrai o backend (FAISS/OpenSearch)
\end{lstlisting}

\subsection{Passo 3 --- Instanciar agente persona}

O modo \texttt{persona} reaproveita contexto do cluster, mas muda voz e
comportamento por meio de \textit{system prompt} dedicado.

\begin{lstlisting}[language=Python,
  caption={Criacao de agente persona por cluster.},
  label={lst:persona_lab}]
agente_persona = _criar_agente_persona(
    cluster_id=cluster_id,
    perfil=perfis.loc[cluster_id],
    retriever=retriever_cluster(cluster_id),
    model=modelo_llm,
)

resposta = agente_persona.invoke({"pergunta": "O que voce espera de um cartao?"})
\end{lstlisting}

\subsection{Passo 4 --- Criar digital twin sob demanda}

No modo \texttt{twin}, o agente nao depende da rotulagem de cluster; ele e
construido com base no registro individual e no retriever filtrado por
\texttt{cliente\_id}.

\begin{lstlisting}[language=Python,
  caption={Twin sob demanda com contexto individual.},
  label={lst:twin_lab}]
twin = criar_twin_sob_demanda(
    cliente_id="C00042",
    dados_cliente=cliente_exemplo,
    model=modelo_llm,
)

out = twin.invoke({"pergunta": "Como posso melhorar meu score de credito?"})
\end{lstlisting}

\subsection{Passo 5 --- Chamar a API em cada modo}

O laboratorio fecha o ciclo testando os tres modos pela API:

\begin{lstlisting}[language=bash,
  caption={Chamadas de validacao dos modos segmento, persona e twin.},
  label={lst:modos_api_lab}]
# segmento
curl -X POST "$API_URL/query" -H "Content-Type: application/json" -d '{
  "modo": "segmento", "pergunta": "Qual oferta para este cliente?", "dados_cliente": {...}
}'

# persona
curl -X POST "$API_URL/query" -H "Content-Type: application/json" -d '{
  "modo": "persona", "cluster_id": 1, "pergunta": "O que voce valoriza em investimentos?"
}'

# twin
curl -X POST "$API_URL/query" -H "Content-Type: application/json" -d '{
  "modo": "twin", "cliente_id": "C00042", "pergunta": "Qual o melhor proximo produto para mim?", "dados_cliente": {...}
}'
\end{lstlisting}

\subsection{Passo 6 --- Resultado completo da aula}

Ao final, o aluno deve conseguir demonstrar:

\begin{itemize}[leftmargin=*]
  \item classificacao automatica do cliente em cluster;
  \item resposta gerencial (segmento), arquetipica (persona) e individual (twin);
  \item coerencia entre contexto recuperado e estilo de resposta por modo.
\end{itemize}

\subsection{Passo 7 --- Evidencias de execucao em producao}

Em teste de smoke test na AWS, executado pelo frontend CloudFront integrado a
API Gateway, os modos \texttt{persona} e \texttt{twin} apresentaram
conclusao com estado \texttt{COMPLETED}, mantendo aderencia de voz em primeira
pessoa e recomendacoes compativeis com o contexto recuperado.

\begin{itemize}[leftmargin=*]
  \item \textbf{Persona} (\texttt{cluster\_id=1}):
        \texttt{request\_id}: \nolinkurl{ddbc5210-63fc-47bc-a274-283e60be84bc}. \par
        Estado final \texttt{COMPLETED};
  \item \textbf{Twin} (\texttt{cliente\_id=C00042}):
        \texttt{request\_id}: \nolinkurl{e5e348e7-d266-4ae9-9f26-eb0e659a56a7}. \par
        Estado final \texttt{COMPLETED}.
\end{itemize}

Como resultado, observa-se que a estrategia de segmentacao e a construcao do
agente individual permanecem operacionalmente consistentes no ambiente
implantado, e nao apenas em execucao local de laboratorio.

% ---------------------------------------------------------------
\section{Plano de testes do capitulo}
% ---------------------------------------------------------------

\subsection{Teste funcional}

\begin{itemize}[leftmargin=*]
  \item \textbf{TC-01:} classificar 50 clientes e verificar ausencia de erro
        de dimensao no \textit{scaler}/K-Means.
  \item \textbf{TC-02:} executar \texttt{modo=persona} com \texttt{cluster\_id}
        explicito e retornar resposta em primeira pessoa.
  \item \textbf{TC-03:} executar \texttt{modo=twin} com \texttt{cliente\_id}
        e confirmar uso de contexto individual.
\end{itemize}

\subsection{Teste de desempenho}

\begin{itemize}[leftmargin=*]
  \item \textbf{TP-01:} tempo de classificacao por cliente $\leq 50$\,ms.
  \item \textbf{TP-02:} latencia mediana da inferencia por modo
        $\leq 5$\,s em \texttt{dev}.
\end{itemize}

\subsection{Teste de qualidade de resposta}

\begin{itemize}[leftmargin=*]
  \item \textbf{TQ-01:} em 10 perguntas por modo, verificar aderencia de tom
        (3a pessoa em segmento; 1a pessoa em persona/twin).
  \item \textbf{TQ-02:} avaliar fidelidade com LLM-as-Judge;
        criterio: media $\geq 8{,}0/10$ por modo.
\end{itemize}

% ---------------------------------------------------------------
\section{Criterios de aceite}
% ---------------------------------------------------------------

\begin{itemize}[leftmargin=*]
  \item TC-01 a TC-03 aprovados.
  \item TP-01 e TP-02 dentro dos limites definidos.
  \item TQ-01 com conformidade de voz em pelo menos 90\% dos casos.
  \item Execucao real em producao com modos \texttt{persona} e \texttt{twin}
        finalizando em \texttt{COMPLETED}.
\end{itemize}

% ---------------------------------------------------------------
\section{Espaco para expansao iterativa}
% ---------------------------------------------------------------

\begin{itemize}[leftmargin=*]
  \item \textbf{v8.1} --- Incorporar cache LRU para twins de alta frequencia.
  \item \textbf{v8.2} --- Evoluir para clusterizacao periodica com janela movel.
  \item \textbf{Lacuna} --- Investigar validacao causal para evitar vies de segmentacao.
\end{itemize}
